\subsection{Internal validity}
\label{sec:limitations-internal}

The thresholds of the predictions were set with the committed values in view: the ratio of one half for the gap index, the bias of $-10\%$ for the treatments, the rejection rate of 10\% for the tests. The new seeds of the extension guard against a dependence on the seed of the committed computation; they do not guard against the choice of thresholds. The interval of the gap index and of the risk-weight difference resamples the annual cohorts as if they were independent, whereas rating statistics have serial dependence through the cycle; a block resample would widen the interval \citep{Kunsch_1989}, and the intervals reported are therefore likely too narrow for the real-data comparison. The Monte Carlo standard errors of the known-truth results are small ($\removedBkNinese$ percent for the bias of B at $k=9$) and cover the sampling of the generator only, not the estimation error of the generator, which is itself fitted to S\&P's counts. The transition tests were run at three levels of $k$ only, and three of the five were run after the first two, once their results were known; the threshold of 10\% for H3 was fixed with the committed rates in view, and the statement that no matrix test sees the withdrawals was tested more narrowly than it was made until the further tests were run. The duration estimator is measured on one generator and one withdrawal window, and the Aalen-Johansen estimator is not measured. The restricted form of the conjecture of Section~\ref{sec:theory-conjecture} is a description of the same measurements, not a test.

\subsection{External validity}
\label{sec:limitations-external}

The known-truth results hold for the generator they were run on: a Markov chain fitted to S\&P's pooled annual counts, cohorts of the size of S\&P's mean cohort, a baseline withdrawal rate of $\wRate\%$ a year and an elevated window of one year before default. A different agency, a different rate or window, or an informative mechanism that acts through the rating path rather than through the time to default would give different biases; the closed forms give the direction and the order of magnitude, and the measurement has to be repeated for each. The real-data results hold for the corporate long-term ratings of the EU entities of three agencies, in the categories scale, for 2000 to 2025; they do not extend to the agencies' global books, to other asset classes or to the notch scale. No data say whether withdrawals of real ratings were informative, and the value $k=9$ is a point on the ladder, not an estimate.

\subsection{Construct validity}
\label{sec:limitations-construct}

Four constructs deserve care. First, the followed treatment is the product's reading of paragraphs 73 and 76 of the guidelines at the one-year level, in which the end state of every rating outstanding at the start is read on the latent path; the guidelines speak of obligors and credit obligations, and the repository counts ratings. Second, the gap between the pages mixes the definition, informative withdrawal and the difference between the cohorts, and the protocol cannot separate them on real data. Third, the risk weight is that of a foundation-approach corporate exposure with a loss given default of $45\%$, a maturity of $2.5$ years, the Basel III function and the PD floor, applied to a portfolio whose grade mix is that of the agency's mean cohort; a bank's portfolio and parameters differ, so the risk-weight differences are an illustration of the size of the effect and not a capital figure. Fourth, the long-run average is the mean of annual rates, as paragraph 84 states, which weights each year equally whatever its cohort; the pooled rates of the same counts differ from it in the speculative grades.

\subsection{Scope}
\label{sec:limitations-scope}

The protocol is a measurement and does not repair the bias it finds. It uses synthetic obligors for the truth and so does not test the treatments on a bank's own data, where the follow-up of withdrawn obligations may be available. It reports the PD of a grade over one year and does not carry the conclusions to multi-year horizons, where the chain of one-year matrices compounds the bias. The thin-grade results assume that the cohort is a set of independent obligors with one PD, which the binomial enumeration requires; with default correlation the coverage would be lower still, and the dependence correction of the effective cohort size, which the engine provides for the Wald and Agresti-Coull intervals, is outside this study. The generator has no cycle: its chain is time-homogeneous, so the study does not measure how withdrawal that is informative only in recessions would change the long-run average, an interaction that the real cohorts may contain and that a generator with a recession regime could reproduce. % voice-ok: withdrawn rating, the domain term of the study
